\documentclass[10pt,a4paper]{amsart}
\usepackage[margin=19mm]{geometry}
\usepackage{amsmath,amssymb,mathtools}
\usepackage[scr=rsfs,cal=euler]{mathalfa}
\usepackage{xfrac}
\usepackage[unicode,hidelinks,bookmarks=false]{hyperref}
\newcommand{\ie}{i.e.\ }
\newcommand{\cf}{cf.\ }
\newcommand{\resp}{respectively\ }
\newcommand{\Subsection}{\subsection}
\providecommand{\nobreakdash}{\nobreak\hbox{-}}
\setlength{\emergencystretch}{2em}
\multlinegap0pt
\usepackage{bm}
\usepackage{bbm}
\usepackage{dsfont}
\usepackage{xspace}
\usepackage{cancel}
\usepackage{mathtools}
\usepackage{amsthm}
\makeatletter
\@ifundefined{theorem}{\newtheorem{theorem}{Theorem}}{}
\@ifundefined{definition}{\newtheorem{definition}[theorem]{Definition}}{}
\@ifundefined{proposition}{\newtheorem{proposition}[theorem]{Proposition}}{}
\makeatother
\newcommand{\R}{\mathbb{R}}
\title[Reilly inequality for Varifolds]{Reilly inequality for Varifolds}
\author{Jean-Fran\c{c}ois Grosjean, Antoine Lemenant, R\'emy Mougenot}
\date{Source: arXiv:2507.08354v1 (CC BY 4.0)}
\begin{document}
\maketitle

\noindent\emph{Source and licence.} Reilly inequality for Varifolds. Version arXiv:2507.08354v1 (CC BY 4.0), distributed under the Creative Commons Attribution 4.0 International licence, as recorded in the arXiv licence metadata for this exact version and shown on the abstract page https://arxiv.org/abs/2507.08354v1. Full rights notice: licenses/arxiv-2507.08354-CC-BY-4.0.txt.\par\smallskip

\noindent\textit{Editorial scope.} Counted unit: the Proposition whose source label is \texttt{Spectrum-Equilateral-Gon} in the article (source0.tex lines 1096--1100, source label \texttt{Spectrum-Equilateral-Gon}), delivered together with the article's own complete original proof of it (source0.tex lines 1102--1147, one \texttt{proof} environment of 46 lines). No mathematics is written, repaired or re-derived: statement and proof are the article's own text line for line. Required context retained uncounted: defPolygon (lines 288--293); lambdaP (lines 677--679); firstorder1 (lines 693--704). Every definition, display and label of those lines is retained unchanged. Omitted from this package and not redistributed: 1--287, 294--676, 680--692, 705--1095, 1101--1101, 1148--1979. Nothing inside a retained excerpt is omitted or altered: every retained body is delivered exactly as the article's lines are written, so no omission receipt of any kind accompanies this package. Citations: the retained excerpts carry no citation command at all, so no bibliography entry of the article is transcribed and no reference is redistributed. Reference adaptation: none was needed and none was made; every reference inside the delivered unit resolves inside it, so no unresolved reference or citation is produced. Printed slips kept literal: no printed word, symbol, hypothesis or number of the counted statement or of its proof was altered. Layout and class: the article's own class is \texttt{smfart} with packages including amssymb, amsfonts, physics, esint, braket, enumitem, appendix, tikz, xcolor, graphicx, epsfig, subfigure, psfrag, listings; the wrapper is the lane's pinned \texttt{amsart} editorial wrapper at 10pt on A4, which restates the theorem-like environments the retained text needs and adds the article's own macro definitions that the retained text needs (\texttt{\textbackslash R}), with extra wrapper packages (bm, bbm, dsfont, xspace, cancel, mathtools). The delivered numbering is the unit's own. Assets: no retained span contains a figure environment, a graphics-inclusion command, a PDF-page inclusion, a table or a TikZ picture; no image file is copied into this package, the package declares no asset dependency and no image file is redistributed with it. Licence: version arXiv:2507.08354v1 of this article is distributed under the Creative Commons Attribution 4.0 International licence, as recorded in the arXiv licence metadata for this exact version and shown on the abstract page https://arxiv.org/abs/2507.08354v1. A full rights notice is delivered at licenses/arxiv-2507.08354-CC-BY-4.0.txt. Characters: every retained excerpt is pure ASCII, so the delivered engine, XeLaTeX with its default Latin Modern text font, reports no missing character.
\medskip\begin{definition} \label{defPolygon} For $n\geqslant 3$ we say that  $A=\{A_i\}_{1\leqslant i \leqslant n+1}$ is  a polygon with vertices $A_i \in \R^N$ if  $A_{n+1}=A_1$ with no self-intersection points, i.e. we assume that  the open  segment $]A_i,A_{i+1}[$ does not intersect any another segment $]A_j,A_{j+1}[$, and we also assume that each vertex $A_i$ belongs to the two (and only two) consecutive segments $[A_{i-1},A_i]$ and $[A_i,A_{i+1}]$. We then define the rectifiable set $K_A$ defined by the union of the segment joining the vertices $A_i$, namely
 $$K_A:=\bigcup_{1\leqslant i \leqslant n}[A_i , A_{i+1}],$$
and the associated rectifiable varifold $V_A$ with constant multiplicity one  defined through 
 $$\int \varphi \;dV_{A}:=\int_{K} \varphi(x,P(x)) \;d\mathcal{H}^1(x).$$
Here $P(x)$ is the approximate tangent plane thus equals $\textnormal{Vect}\{A_{i+1}-A_{i}\}$ on the segment $[A_i,A_{i+1}]$. 
\end{definition}

 \begin{eqnarray}
 \lambda_1^A:=\inf_{u \in C^\infty(\mathbb{R}^N) \setminus \{0\} \text{ s.t. } \int_{\R^N} u \; d \mu_A=0} \frac{\int |\nabla_P u(x)|^2 \;d V_A(x,P)}{\int_{\mathbb{R}^N} u^2\;d\mu_A}. \label{lambdaP}
 \end{eqnarray}

 \medskip\begin{proposition}[First order conditions] \label{firstorder1}  The infimum in \eqref{lambdaP} is achieved  in the space $H^1(K_A)$ for a function $u$ that satisfies the following first order condition: let  $T_i$ be the parameterization of $[A_i,A_{i+1}]$ defined by $T_i(t) := (1-t)A_i + tA_{i+1}$. We also denote by $\ell_i:=|A_{i+1}-A_i|$. Then there are  some constants $\alpha_i,\beta_i \in \R$ such that
$$
\left\{
 \begin{array}{ll}
u\circ T_i(t)=\alpha_i +t \beta_i, \quad  &\textnormal{ for } t\in[0,1] \textnormal{ and }  1\leqslant i \leqslant n,\\
\alpha_i+\beta_i=\alpha_{i+1},\\\displaystyle 
\lambda_1\alpha_i=\left(\frac{\beta_{i-1}}{\ell_{i-1}} - \frac{\beta_i}{\ell_i}\right), \quad &\textnormal{for } 1\leqslant i  \leqslant n.
 \end{array}
  \right. 
 $$
and $\alpha_{n+1}=\alpha_1$ and $\beta_{n+1}=\beta_1$.
 \end{proposition}

\medskip\begin{proposition}\label{Spectrum-Equilateral-Gon} Let $A=\{A_i\}_{1\leqslant i \leqslant n+1}$ be a polygon in $\R^N$ as in Definition \ref{defPolygon}. Assume that $A$ is equilateral and $L$ is the perimeter of $A$. Then  the spectrum of $A$ is 
$$\textnormal{Spec}(A)=\left\{\frac{4n}{L}\sin^2\frac{k\pi}{n}\ ; \ k\in\left\{0,\cdots\lfloor\dfrac{n}{2}\rfloor\right\}\right\}.$$
In particular,
$$\lambda_1^A=\frac{4n}{L}\sin^2\left(\frac{\pi}{n}\right).$$
\end{proposition}

\begin{proof} Since $A$ is equilateral, for any $i\in\{1,\cdots,n\}$, $|A_{i+1}-A_i|=\ell=L/n$. Let $u$ be an eigenfunction for an eigenvalue $\lambda$. From Proposition \ref{firstorder1} we know that if  $T_i$ is the parameterization of $[A_i,A_{i+1}]$ defined by $T_i(t) := (1-t)A_i + tA_{i+1}$, then there exists some constants $\alpha_i,\beta_i \in \R$ such that $u\circ T_i(t)=\alpha_i +t \beta_i$ for $t\in[0,1]$ and $ 1\leqslant i \leqslant n$, and 
$$
\left\{
 \begin{array}{ll}
\alpha_i+\beta_i=\alpha_{i+1},\\
\lambda\alpha_{i+1}=\dfrac{1}{\ell}(\beta_i-\beta_{i+1}), \quad \textnormal{ for } 1\leqslant i  \leqslant n,
 \end{array}
  \right. 
 $$
with $\alpha_{n+1}=\alpha_1$ and $\beta_{n+1}=\beta_1$. Note that this system above is equivalent to
\begin{equation}\label{equivalence-system}
 \begin{pmatrix}
\alpha_{i+1}\\\beta_{i+1}
\end{pmatrix}=M\begin{pmatrix}
\alpha_i\\\beta_i
\end{pmatrix}:=\begin{pmatrix}
1&1\\-\Lambda&1-\Lambda
\end{pmatrix}\begin{pmatrix}
\alpha_i\\\beta_i
\end{pmatrix},\quad \textnormal{ for }1\leqslant i  \leqslant n,
\end{equation}
with $\Lambda = \lambda \ell$. Moreover, $$M^n\begin{pmatrix}
\alpha_1\\\beta_1
\end{pmatrix}=\begin{pmatrix}
\alpha_1\\\beta_1
\end{pmatrix}.$$ Then $1$ is an eigenvalue of $M^n$ and it follows that the spectrum of $M$ contains an $n$-th root of $1$. A straightforward computation shows that the characteristic polynom of $M$ is $P=(X-1)^2+\Lambda(X-1)+\Lambda$ and the discriminant of $Y^2+\Lambda Y+\Lambda$ is $\Delta=\Lambda(\Lambda-4)$. Look for the conditions on $\Lambda$ so that the spectrum of $M$ contains an $n$-th root of $1$. 

\emph{Case} $\Lambda=0$. Imediatly, $\lambda=0$ and $\beta_i$ are constant. Then $\alpha_{n+1}=\alpha_1+n\beta_1=\alpha_1$, $\beta_i=0$ and $\alpha_i$ are constant. It follows that $u$ is constant. Conversely, a constant function is an eigenfunction for $\lambda=0$. 

 \emph{Case} $0<\Lambda<4$. One has $\Delta<0$ and the roots of $P$ are $x^{\pm}=(2-\Lambda)/2\pm i\sqrt{\Lambda(4-\Lambda)}/2$. It is easy to prove that $|x^{\pm}|=1$ and $x^+$ or $x^-=\overline{x^+}$ is an $n$-th root of $1$ if and only if $x^+$ or $\overline{x^+}=\exp(i2k\pi/n)$ for $1 \leqslant k \leqslant n-1$, in fact $k\neq0$ since $\Lambda\neq0$. This is  equivalent to say that $(2-\Lambda)/2=\cos(2k\pi/n)$ and $\Lambda=4\sin^2(k\pi/n)$.  From this we obtain that $$\lambda=\dfrac{4n}{L}\sin^2\dfrac{k\pi}{n}.$$ 
 
\emph{Case }$\Lambda=4$. We have $x=-1$ and $x$ is an $n$-th root of $1$ if and only if $n$ is even. In this case $$\lambda=\dfrac{4n}{L}=\dfrac{4n}{L}\sin^2\dfrac{(n/2)\pi}{n}.$$

\emph{Case } $\Lambda>4$. The roots $x^{\pm}\in\R$. If $x^{\pm}$ is an $n$-th root of $1$ then $x^{\pm}=\pm 1$ which is impossible. Thus  we have proved that the spectrum of $A$ is including in 
$$\textnormal{Spec}(A)\subset \left\{\frac{4n}{L}\sin^2\frac{k\pi}{n}\ \vert \ 0\leqslant k \leqslant n-1\right\}=\left\{\frac{4n}{L}\sin^2\frac{k\pi}{n}\ \vert\ 0 \leqslant k \leqslant \lfloor\dfrac{n}{2}\rfloor\right\}.$$
Conversely, let $k\in\{1,\cdots,n-1\}$, $\lambda=(4n/L)\sin^2(k\pi/n)$ and $\Lambda=\lambda\ell$ (the case $k=0$ corresponds to $\lambda=0$ already discussed above). Then $\exp(i2k\pi/n)$ is an eigenvalue of $M$ and $1$ is an eigenvalue of $M^n$. Let $(
\alpha_1,\beta_1)$ be an eigenvector of $M^n$ and for any $i\in\{2,\cdots,n+1\}$ define $$\begin{pmatrix}
\alpha_i\\\beta_i
\end{pmatrix}=M^{i-1}\begin{pmatrix}
\alpha_1\\\beta_1
\end{pmatrix}.$$ Then for any $i\in\{1,\cdots,n\}$, $$ \begin{pmatrix}
\alpha_{i+1}\\\beta_{i+1}
\end{pmatrix}=M\begin{pmatrix}
\alpha_i\\\beta_i
\end{pmatrix}.$$ If we define a function $u$ as above with scalar $\alpha_i$ and $\beta_i$, the equivalence \ref{equivalence-system} allows us to say $u$ is an eigenfunction for $\lambda=\frac{4n}{L}\sin^2\frac{k\pi}{n}$. Consequently, $$\textnormal{Spec}(A)=\left\{\frac{4n}{L}\sin^2\frac{k\pi}{n}\ ; \ k\in\left\{0,\cdots\lfloor\dfrac{n}{2}\rfloor\right\}\right\}.$$
\end{proof}  

\end{document}
